Un pendule est constitué d'une petite boule métallique, de masse
$m=\qty{80}{\g}$, suspendue à un fil inextensible, de masse considérée comme
négligeable et de longueur $\ell=\qty{1.00}{\m}$.

Le fil est accroché en un point fixe $O$ et les mouvements du pendule
s'effectuent dans un plan vertical. Le fil du pendule étant initialement
vertical, on l'écarte de cette position d'un angle $\theta_{m}=\ang{45}$.

Puis, fil tendu, on le lâche sans vitesse (position~1).

\begin{figure}[H]
    \centering
    \includegraphics[width=.5\textwidth]{0ef5c135-036a-4cbb-904f-bb06b1b6e281-5}
\end{figure}


\begin{enumerate}
    \item Décrire le mouvement de la boule.

          Justifier la conservation de la somme $E_{c}+E_{pp}$ pour la boule du
          pendule.
    
          Quelle est la transformation d'énergie qui s'effectue au cours du
          mouvement~?
    \item Déterminer la valeur $v_{2}$ de la vitesse de la boule lorsqu'elle
          passe par la position verticale (position~2).
    
          Justifier le calcul.
    
          Faire le calcul littéral et procéder à l'application numérique.
    \item La position intermédiaire du pendule est définie par l'angle $\theta$
          qu'il forme avec la verticale; la valeur de la vitesse de la boule est
          alors $v$.
        
          On fait l'hypothèse que l'énergie potentielle de pesanteur est nulle
          dans la position la plus basse que le pendule peut occuper
          (position~2).

          Écrire, avec cette convention, l'équation de conservation régissant le
          mouvement du pendule.
    
          En déduire la formule littérale donnant la valeur $v$ de la vitesse en
          fonction de $\theta$, $\theta_{m}$, $g$ et $\ell$. Faire l'application
          numérique pour $\theta=\ang{30}$ et $\theta^{\prime}=\ang{15}$.
    
    \item On place, sur la verticale de $O$, à la distance $d=\qty{60.0}{\cm}$,
          une tige métallique $T$ sur laquelle le fil du pendule, lâché comme
          précédemment ($\theta_{m}=\ang{45}$), vient buter.
    
          \begin{figure}[H]
              \centering
              \includegraphics[width=.5\textwidth]{0ef5c135-036a-4cbb-904f-bb06b1b6e281-5(1)}
          \end{figure}
          
          En admettant que le choc du fil du pendule sur la tige $T$ s'effectue
          sans perte d'énergie, déterminer l'angle $\alpha$ dont le pendule
          remonte après le choc avec la tige.
\end{enumerate}
